\cleardoubleoddpage
\chapter{Results}
\label{chap:results}

\section{Evidence and Selection Integrity}

The accepted release is generated from the frozen protocol tag
\texttt{protocol/multiseed-v1} (source commit prefix \texttt{be976808}; the full digest is in
Appendix~\ref{app:reproducibility}). The
dataset revision is the byte-verified consumed universe with revision digest
\texttt{cyclone-consumed-bytes-sha256:ff2867e9}; the universe and outer-fold manifest digests are
\texttt{96c85a70} and \texttt{da5d95b8}, respectively (the complete values are in the release
manifest). The evidence tree contains 255 planned stage slots: 230 accepted ledger slots (225
metric stages and five selection barriers) and 25 skipped because
their learned family was not selected. There are no invalidated accepted rows.

Every accepted metric stage is finite and stable, uses the declared fold and training seed, and records
the expected checkpoint/cache lineage. The five selection records all report
\code{test\_evidence\_opened=false}. The W\&B contract is disabled and verified for every accepted metric stage.
These checks matter because a numerically plausible rollout can still be scientifically invalid if
an encoder has seen its supposed test trajectories or if an unselected test model was inspected.

\section{Qualitative Latent-Space Geometry}

\autoref{fig:latent-space} visualizes all 1,173 encoded snapshots from one reconstructed
protocol run: outer fold 0, training seed 52. The trajectory roles come directly from the frozen
fold manifest, rather than from a second random split. Colour represents the preprocessed scalar
flux target; marker shape identifies the training, validation, or test role. PCA gives a linear
view, while t-SNE at perplexities 5 and 30 gives two local-neighbourhood views under the same seed.

\begin{figure}[htbp]
  \centering
  \includegraphics[width=\textwidth]{latent_space_projection.png}
  \caption{Qualitative latent-space geometry for all 1,173 points in the canonical
  \texttt{multiseed-v1} outer-fold-0, seed-52 cache. Colour denotes preprocessed flux and marker
  shape denotes the exact manifest role. This is one protocol-bound run, not an aggregate
  multiseed performance estimate. The t-SNE axes and global distances have no physical meaning,
  and the panels do not establish split separation or model superiority.}
  \label{fig:latent-space}
\end{figure}

The projections are descriptive only. There is no shared latent coordinate system across the 25
fold--seed encoder fits, so averaging their coordinates would not define a meaningful multiseed
plot. The visible flux ordering is compatible with the supervised diagnostic objective, but it is
not a held-out accuracy result and does not imply that the representation retains the full
five-dimensional state.

\section{Outer-Fold Selection}

The fold-level family choices are shown in \autoref{tab:fold-selection}. The selection means are
computed from five matched validation seeds; the table does not use outer-test metrics. The
variation in selected family is itself a result: the evidence does not justify naming a single
globally superior architecture.

\begin{table}[htbp]
  \centering
  \caption{Architecture selected independently in each outer fold. Selection uses the mean
  validation trajectory-balanced flux RMSE over seeds 52--56; lower is better.}
  \label{tab:fold-selection}
  \small
  \begin{tabular}{lrrr}
    \toprule
    Outer fold & Test trajectories & Selected family & Validation selection rule \\
    \midrule
    0 & 11 & Transformer & five-seed mean flux RMSE \\
    1 & 10 & Transformer & five-seed mean flux RMSE \\
    2 & 10 & GRU & five-seed mean flux RMSE \\
    3 & 10 & Transformer & five-seed mean flux RMSE \\
    4 & 10 & GRU & five-seed mean flux RMSE \\
    \bottomrule
  \end{tabular}
\end{table}

The selection is deliberately two-level. Within each seed, a sequence checkpoint is chosen by
validation latent RMSE. The family comparison then averages the resulting validation flux RMSE
over all five training seeds. This distinction prevents the flux metric from silently becoming a
test-aware checkpoint rule and prevents a single lucky initialization from determining the family.

\section{Primary Diagnostic Result}

The primary estimand is the selected learned model minus direct observed-flux persistence in
trajectory-balanced flux RMSE. The aggregate learned error is 14.6729, whereas observed
persistence is 10.0207. The difference is therefore +4.6522 in favour of the baseline. The paired
hierarchical bootstrap interval is [+2.5788, +6.6470], based on 10,000 replicates and the frozen
bootstrap seed. It is entirely above zero, so the protocol's negative-result rule applies: the
learned surrogate shows no convincing advantage and is, on this metric and universe, worse than
the simple observed diagnostic persistence.

\begin{table}[htbp]
  \centering
  \caption{Primary outer-test estimate across five folds and five matched seeds. Positive
  differences mean that the learned model has larger error than observed persistence.}
  \label{tab:primary-results}
  \small
  \begin{tabularx}{\textwidth}{>{\raggedright\arraybackslash}X r >{\raggedright\arraybackslash}X}
    \toprule
    Quantity & Value & Interpretation \\
    \midrule
    Selected learned flux RMSE & 14.6729 & lower is better \\
    Observed-flux persistence flux RMSE & 10.0207 & primary baseline \\
    Learned minus observed difference & +4.6522 & baseline advantage \\
    Hierarchical bootstrap 95\% interval & [+2.5788, +6.6470] & excludes zero \\
    Hierarchically weighted improved fraction & 31.6\% & learned error lower \\
    Paired fold--seed--trajectory rows & 255 & finite and paired \\
    \bottomrule
  \end{tabularx}
\end{table}

The hierarchically weighted fraction is not a claim that exactly 31.6\% of raw rows are an
independent sample; it gives equal weight to folds and seeds before trajectory averaging. The
per-fold point differences are positive in all five folds: +1.4189, +3.9495, +6.5372, +7.2051,
and +4.1503. Thus the negative conclusion is not caused by one pathological outer split, although
the magnitude varies substantially between folds.

\begin{figure}[htbp]
  \centering
  \includegraphics[width=0.78\textwidth]{verified_primary_by_fold.png}
  \caption{Selected-model minus observed-persistence flux RMSE by outer fold. Points are fold
  means and bars are the release's fold-level bootstrap intervals. The zero line marks equal
  performance; every fold is on the positive-error side.}
  \label{fig:primary-by-fold}
\end{figure}

\section{Latent Dynamics and Forecast Horizon}

Decoded latent-state persistence has flux RMSE 14.8049, only slightly worse than the selected
learned models at the aggregate level. This is a useful reference because it asks whether a
sequence model improves on holding the encoded state fixed, but it is not the primary diagnostic
baseline: it still passes through the learned diagnostic head. The diagnostic-head oracle, which
receives the true future latents, reaches flux RMSE 12.3732. It is better than either latent
rollout but remains worse than observed persistence. The result therefore identifies a material
representation-to-diagnostic gap without isolating a single cause; the oracle is a control, not a
forecast ceiling.

The secondary latent estimand is selected learned minus latent-persistence latent MSE. Learned
latent MSE is 0.4497 versus 0.4307 for persistence, giving +0.0189. Its interval
[-0.0699, +0.1332] includes zero. Accordingly, the experiment does not support a claim that the
sequence model improves latent dynamics either. The point estimate is compatible with a small
regression or improvement under another realization, and the interval is the correct summary of
that uncertainty.

\begin{figure}[htbp]
  \centering
  \includegraphics[width=0.78\textwidth]{verified_latent_mse_by_step.png}
  \caption{Mean latent MSE by forecast step across outer folds. The ribbon is the standard
  deviation between outer-fold means; it is not a confidence interval. Curves are generated from
  the release's fold-level by-step arrays rather than from a selected seed.}
  \label{fig:latent-horizon}
\end{figure}

\begin{figure}[htbp]
  \centering
  \includegraphics[width=0.78\textwidth]{verified_flux_rmse_by_step.png}
  \caption{Mean diagnostic flux RMSE by forecast step for the selected learned model, observed
  persistence, and the diagnostic-head oracle. Shaded ribbons show between-outer-fold standard
  deviation, not confidence intervals. The figure is a horizon diagnostic, not a new selection
  criterion.}
  \label{fig:flux-horizon}
\end{figure}

The horizon curves remain finite, but finiteness is a stability check rather than an accuracy
claim. In particular, a rollout can be numerically stable while drifting systematically away from
the observed diagnostic. The positive aggregate difference and the per-fold results are therefore
more important than the absence of exploding values.

\section{Spectral Diagnostics}

Spectral results are kept per target, as required by the protocol. For $\code{fluxspec}$, the
learned relative $L_2$ error is 1.5026 with shape correlation 0.6011, compared with 0.4433 and
0.9123 for observed persistence. For $\code{kyspec}$, the learned relative $L_2$ error is 28.1003
with shape correlation 0.7226, compared with 0.5296 and 0.9330 for observed persistence. Decoded
latent persistence is similarly worse than observed persistence (relative $L_2$ 1.4597 for
\code{fluxspec} and 27.2029 for \code{kyspec}).

These large, target-dependent relative errors are not combined into a single claim of spectral
fidelity. The shape correlations show that some broad spectral ordering is retained, while the
scale-sensitive errors show that this is insufficient for quantitative prediction. Spectra
therefore reinforce, rather than overturn, the primary flux result.

\section{Answer to the Research Questions}

RQ1 asks whether the compact state retains diagnostic information. The finite oracle and nontrivial
spectral correlations show that the representation carries usable signal, but the oracle's error
relative to observed persistence shows that it is not sufficient for accurate diagnostic rollout.
RQ2 asks whether a learned sequence model beats observed persistence. Across the frozen nested
cross-validation estimate, it does not: the interval for learned minus observed flux RMSE is
strictly positive. RQ3 asks whether the cache-normalized architecture is a reliable winner. The
fold selections alternate between GRU and Transformer, so no global winner is supported.

No historical seed-52 values, mixed-split runs, unpaired exploratory baselines, or external
architecture numbers are used in these conclusions. The negative result is bounded to the
51-trajectory universe, the preprocessing contract, the eight-step horizon, and the declared
diagnostic targets.
